\section{Related Work}
\label{sec:app:lit}

\textbf{Hyper-parameter optimization.} There is significant prior work in hyper-parameter optimization \cite{agrawal2019dodge, cowen2022hebo, li2017hyperband, bergstra2011algorithms, bergstra2012random, falkner2018bohb, eriksson2019scalable, ansel:pact:2014}. Indeed, as learning systems become more intricate, it is crucial that we eke out the most performance. However, this is a non-trivial problem, as evidenced by the long line of research in this direction. 

The simplest form of hyper-parameter search is random search, which tries $n$ randomly chosen hyper-parameter configurations. Opentuner \cite{ansel2014opentuner} is a multi-armed bandit meta-technique with a sliding window that incorporates an exploration/exploitation trade-off based on the number of times a specific technique is used. It combines DE, a greedy bandit mutation technique, and hill-climbing methods. 

Bayesian Optimization (BO) has emerged as the most popular technique for HPO. \citet{bergstra2011algorithms} propose the Tree of Parzen Estimators (TPE) algorithm. Rather than model $p(y|x)$, TPE models $p(x|y)$ as 
$$
p(x|y) = \begin{cases}
    l(x) &  y < y^* \\
    g(x) &  y \geq y^*
    \end{cases}
$$
where $y^*$ is chosen so that $p(y < y^*) = \gamma$ for some quantile $\gamma$. The functions $l(x)$ and $g(x)$ are kernel density estimates. TPE optimizes the EI, which they show is equivalent to maximizing $l(x) / g(x)$.
\citet{snoek2012practical} use Gaussian Process (GP) models as the surrogate function in Bayesian optimization. They use Expected Improvement (EI) as the acquisition function. Similarly, \citet{hernandez2014predictive} use predictive entropy search (PES) as the acquisition function. \citet{swersky2014freeze} exploit iterative training procedures in their Bayesian optimization framework, which they call freeze-thaw Bayesian optimization. \citet{snoek2015scalable} use neural networks for modeling distributions over functions that yields an approach that scales linearly over data size (rather than cubically as in GP-based Bayesian optimization). BOHB \cite{falkner2018bohb} combines the BO-based TPE with HyperBand \cite{li2017hyperband}, replacing the initial random configurations with a model-based search. Notably, BOHB uses a single multi-dimensional KDE instead of hierarchical single-dimensional KDEs used by TPE. The authors of HEBO \cite{cowen2022hebo} note that (i) even simple HPO problems can be non-stationary and heteroscedastic (ii) different acquisition functions can conflict. To tackle the former, they use the Box-Cox \cite{box1964analysis} and Yeo-Johnson \cite{yeo2000new} output transformations and the Kumaraswamy \cite{kumaraswamy1980generalized} input transformation. It also uses NSGA-II to optimize a multi-objective acquisition function. TuRBO \cite{dou2023turbo} assumes the hyper-parameter to performance mapping is Lipschitz, and generates pseudo-points to improve convergence of vanilla BO. It also uses an ensemble of learners to predict performance, and if the prediction is poor, uses it to instead update the GP model. Finally, we mention PriorBand \cite{mallik2023priorband}, which incorporates an expert's prior beliefs about good configurations, but maintains good performance even if that prior is bad.

We defer to \cite{feurer2019hyperparameter} and \cite{bischl2023hyperparameter} for recent reviews on hyper-parameter optimization techniques. There is also a long line of work studying neural architecture search, which aims to find optimal architectures for a dataset. We refer the reader to \citet{white2023neural} for a comprehensive review of the field.

% There is also a long (and growing) line of work studying \textit{neural architecture search} (\cite{baker2016designing,elsken2017simple,liu2017hierarchical,liu2018darts,stanley2002evolving,zoph2018learning,xie2018snas,real2019regularized}). This field studies systematic approaches to find optimal neural network architectures for a dataset. For example, \cite{liu2017hierarchical} start with small building blocks and work up to a hierarchical, more complex architecture. \cite{liu2018darts} frame the problem in a way that allows them to use gradient descent to find an optimal architecture. We do not use neural architecture search in this work, but mention it due to its relevance to hyper-parameter optimization; indeed, neural architecture search focuses on one specific subset of hyper-parameters--the network architecture--without focusing on others such as batch size, optimizer-specific hyper-parameters such as learning rate, preprocessing algorithms used, etc. Moreover, neural architecture search approaches can be slow and take days on a modern GPU to find an optimal architecture, for one dataset and for one run. Our methods use statistical analysis to check that our approach is better by a statistically significant margin, and therefore neural architecture search could not be incorporated in our overall framework. Nevertheless, we direct the reader to \cite{elsken2019neural} for a recent literature review of neural architecture search approaches.

Several benchmarks have been proposed for hyper-parameter optimization: notable ones include YAHPO Gym \cite{pfisterer2022yahpo}, HPO-B \cite{arango2021hpo}, and HPOBench \cite{eggensperger2021hpobench}.

\textbf{Flat minima and generalization.} The idea of flat minima was first studied by \citet{hochreiter1994simplifying}. In particular, they define ``flat minima'' as large connected regions where the weights are $\epsilon-$optimal. \citet{hochreiter1997flat} intuit that because sharper minima require higher precision, flatter minima require less bits to describe. They use this intuition to show that flat minima correspond to minimizing the number of bits required to describe the weights of a neural network. This notion of minimum description length (MDL) \cite{rissanen1983universal, grunwald2007minimum} was also exploited early on by \citet{hinton1993keeping}. Flat minima were revisited by \citet{chaudhari2019entropy}, who noted that minima with low generalization error have a large proportion of their eigenvalues close to zero. They then construct a modified Gibbs distribution corresponding to an energy landscape $f$, and minimize the negative local entropy of this modified distribution, and approximate the gradient via stochastic gradient Langevin dynamics (SGLD) \cite{welling2011bayesian}. However, their assumptions were, admittedly unrealistic. \citet{keskar2016large} show that when using large batch sizes, optimizers converge to sharp minima, which are characterized by many large positive eigenvalues of the Hessian. Further, they define the notion of sharpness as a generalization measure as the robustness to adversarial perturbations in the parameter space:
\begin{equation}
\zeta(\boldsymbol w; \epsilon) = \frac{\max\limits_{\abs{\boldsymbol v} \leq \epsilon(\abs{\boldsymbol w} + 1)} f(\boldsymbol w + \boldsymbol v) - f(\boldsymbol w)}{1 + f(\boldsymbol w)}
\label{eq:sharpness}
\end{equation}
and compute this using 10 iterations of L-BFGS-B \cite{byrd1995limited} with $\epsilon = \{ 10^{-3}, 5 \cdot 10^{-4} \}$. In particular, the above is closely related to the largest eigenvalue of $\nabla^2 f(\boldsymbol w)$. This notion of sharpness was also endorsed by \citet{jiang2019fantastic}, who performed a large-scale study of many complexity measures on two datasets, with 2,187 convolutional networks. In particular, they endorse the following metrics for generalization: (i) variance of gradients (ii) squared ratio of magnitude of parameters to magnitude of perturbation, à la \citet{keskar2016large} (iii) path norm \cite{neyshabur2017exploring} (iv) VC-dimension (inversely correlated).
    
A line of work in physics \cite{baldassi2015subdominant,baldassi2016unreasonable} showed that in the discrete weight scenario (a much more difficult problem), isolated minima were rare, but there existed accessible, dense regions of subdominant minima, and that these were robust to perturbations and generalized better. These authors devised algorithms explicitly designed to search for nonisolated minima. In the continuous weight space, nonisolated minima correspond to flat minima. \citet{dziugaite2017computing} obtain nonvacuous generalization bounds for deep overparameterized neural networks using the PAC-Bayes framework \cite{mcallester1999pac}. \citet{neyshabur2014search} showed that increasing the number of hidden units (which in turn, increases the number of trainable parameters) can lead to a decrease in generalization error with the same training error. \citet{neyshabur2017exploring} showed that sharpness as computed by \eqref{eq:sharpness} is not sufficient to capture the generalization behavior (but noting that ``combined with the norm, sharpness does seem to provide a capacity measure''), and advocate for expected sharpness in the PAC-Bayesian framework, similar to \citet{dziugaite2017computing}. They show that plots of expected sharpness versus KL divergence in PAC-Bayes bounds for varying dataset sizes capture generalization well. \citet{li2018visualizing} showed that the sharpness of the loss surface correlates well with the generalization error. In their seminal paper, \citet{jastrzkebski2017three} showed that SGD is a Euler-Maruyama discretization of a stochastic differential equation whose dynamics are influenced by the ratio of learning rate to batch size (which they call ``stochastic noise''), and that SGD finds wider minima with higher stochastic noise levels than sharper minima. \citet{wu2023implicit} study the flat minima hypothesis through the lens of dynamical stability, and show that SGD will escape from overly sharp (measured by the Frobenius norm of the Hessian), low-loss areas exponentially fast. We note that they use the associate empirical Fisher matrix (AEFM) as an approximation for the Hessian, which holds for low empirical risk (and converges to the Hessian, see \citet{kunstner2019limitations}).

In search of flatter loss surfaces, \citet{seong2018towards} propose the use of non-monotonic learning rate schedules. They advocate for large learning rates, which enable the optimization algorithm to escape sharp minima, and descend into flatter minima. The seminal works of \citet{dauphin2014identifying} and \citet{choromanska2015loss} showed both theoretically and empirically that local minima are more likely to be located close to the global minimum. In particular, \citet{dauphin2014identifying} showed using the perspectives of random matrix theory (via the eigenvalue distribution of Gaussian random matrices \cite{wigner1958distribution}), statistical physics (via the analysis of critical points in Gaussian fields by \citet{bray2007statistics}), and neural network theory \cite{saxe2013exact} that saddle points are exponentially less likely than local minima.

Most recently, \citet{wen2023sharpness} showed that sharpness is neither necessary, nor sufficient for generalization, by studying some simple architectures, showing that generalization depends on the data distribution as well as the architecture; for example, merely adding a bias to a 2-layer MLP makes generalization impossible for the XOR dataset.
    
Of course, it is not possible to discuss generalization in deep learning without discussing the results of \citet{zhang2021understanding}, who showed that deep learners can fit with zero training error on random labels using an architecture that generalizes well when fit to the correct labels. \citet{bartlett2019nearly} and \citet{harvey2017nearly} found that the VC-dimension for deep ReLU networks is $\mathcal{O}(WL \log W)$ and $\Theta(WU)$ respectively, where $W$ is the number of parameters, $L$ the number of layers, and $U$ the number of units.  \citet{hardt2016train} show that stochastic gradient methods are uniformly stable\footnote{An algorithm $A$ is $\epsilon-$uniformly stable if $\forall S, S^\prime \in Z$ for some space $Z$, such that the datasets $S$ and $S^\prime$ differ by at most one example, $\sup\limits_z \mathbb{E}_A [f(A(S); z) - f(A(S^\prime); z)] \leq \epsilon$}, which implies generalization \textit{in expectation}.